\scanpage{22}

\deleted{\unsure{Lemme 3}} Alors on a $x_N=x_{N+1}$.

\uuline{Dém} Posons
\[
\begin{cases}
\deleted{\unsure{y_0=x_{N+1}-x_N}}\\
y_\varepsilon=x_{N+p^\varepsilon}-\unsure{x_N}\qquad0\leq\varepsilon\leq e.
\end{cases}
\]
d'où
\[
y_\varepsilon-y_0=x_{N+p^\varepsilon}-x_{N+1}
\]
\deleted{$\unsure{y'_\varepsilon=x_{N+p^\varepsilon}-x_{N+1}}$}

L'hypothèse s'écrit donc, en distinguant les cas $i=N$, $i=N+1$
\deleted{$\unsure{(N+1)y_0=0}$}\qquad\deleted{$\unsure{(i=N,\ j=N+p^\varepsilon)}$}
\[
(\unsure{S|\overline S})=S(N,e,(y_\varepsilon)_{0\leq\varepsilon\leq e})\quad
\begin{cases}
c_{N,p^\varepsilon}y_\varepsilon=0&0\leq\varepsilon\leq e\\
c_{N+1,p^\varepsilon-1}(y_\varepsilon-y_0)=0&0\leq\varepsilon\leq e\quad(\text{le cas }\varepsilon=0\text{ étant trivial})
\end{cases}
\]
La conclusion s'écrit $y_0=0$. En fait, nous allons déduire cette conclusion du système plus faible \inserted{(équivalent si $E$ est annulé par $p$, i.e. est un $\mathbb F_p$-vectoriel)}
\[
(\overline S')=\overline S(N,e,(y_\varepsilon)_{0\leq\varepsilon\leq e})\quad
\begin{cases}
y_\varepsilon=0\text{ si }c_{N,p^\varepsilon}\not\equiv0\ (p)&0\leq\varepsilon\leq e\\
y_\varepsilon-y_0=0\text{ si }c_{N+1,p^\varepsilon-1}\not\equiv0\ (p)&\unsure{1|0}\leq\varepsilon\leq e.
\end{cases}
\]

Pour prouver que ce système implique $y_0=0$, on récurre sur $N$. Comme la première équation pour $\varepsilon=0$ donne
\[
y_0=0\text{ si }N+1\not\equiv0\ (p)
\]
la conclusion est triviale pour $N=0$. \deleted{\unsure{Alors on}} On suppose donc $N>0$, \inserted{et prouvé pour les $N'<N$.} La conclusion est triviale si $N+1\not\equiv0\ (p)$.
Donc OPS
\[
N+1\equiv0\ (p)\quad\text{soit}\quad N+1=(N'+1)p
\]
où
\[
N'=\frac{N+1}{p}-1<N\qquad\left(\text{car }\frac{N+1}{p}<N+1\right).
\]
Alors le cor. précédent montre que le syst. $\overline S=\overline S(N,e,(y_\varepsilon)_{0\leq\varepsilon\leq e})$ \inserted{(compte tenu de $N+1\equiv0\ (p)$ qui nous permet d'oublier la première équation pour $\varepsilon=0$)} équivaut au système (où $e'=e-1$, $y'_\varepsilon=y_{\varepsilon+1}$ pour $0\leq\varepsilon\leq e'=e-1$)
\[
\deleted{\overline S'=\overline S'(N',e',(y'_\varepsilon)_{\unsure{1\leq\varepsilon}})}\quad
\begin{cases}
y'_\varepsilon=0\text{ si }c_{N',p^\varepsilon}\not\equiv0\ (p)&0\leq\varepsilon\unsure{=|\leq}e'\\
y'_\varepsilon-y_0=0\text{ si }c_{N'+1,p^\varepsilon-1}\not\equiv0\ (p)&\unsure{0|1}\leq\varepsilon\leq e'
\end{cases}
\]
Pour $\varepsilon=0$ cette dernière équation donne \deleted{$\unsure{x}$}
\[
y'_0=y_0
\]
donc on peut se borner à $1\leq\varepsilon\leq e'$ (et se rappeler $y'_0=y_0$).

\mspage{21}
